\section{MULTIMODAL LOCOMOTION AND CONTROL}

\subsection{Walking}

Walking uses alternating tripods. One tripod supports the body while the other is
lifted and advanced; the roles are then exchanged. The final paper must report
joint trajectories in radians rather than raw Dynamixel counts, together with
the policy or state machine that sets phase, stride, and body height.

\subsection{Quasi-Rolling}

For flat-ground transport, the six arc-shaped distal links are aligned into two
longitudinal groups and commanded in velocity mode. Their staggered contact
phases approximate continuous rolling support even though each distal link is an
open arc rather than a complete wheel. Direction changes are produced by
differential left-right velocity commands.

\subsection{Stair Climbing}

During stair ascent, the arcs engage successive stair edges while proximal joints
adjust the front and rear body height. The controller permits a phase offset
between front and rear legs so the robot can accommodate tread spacing that does
not exactly match the nominal leg spacing.

\subsection{Mode Transitions}

\begin{figure}[t]
  \centering
  \fbox{\parbox[c][27mm][c]{0.92\columnwidth}{\centering
    Final Fig.~3: state diagram for walk $\leftrightarrow$ roll
    $\leftrightarrow$ climb, including transition guards and times.}}
  \caption{Mode-transition controller. Each transition is executed without
  mechanical reconfiguration or external handling.}
  \label{fig:mode-transition}
\end{figure}

\draftnote{Document the actual control frequency, interpolation, velocity limits,
current/torque limits, and emergency stop. If reinforcement learning is not used
in the reported experiments, do not add it merely as a future-looking method.}
